\documentclass[conference]{IEEEtran}
\usepackage{cite}
\usepackage{amsmath}
\usepackage{graphicx}
\usepackage{booktabs}
\usepackage{url}
\graphicspath{{../../results/figures/}}

\begin{document}

\title{Next-Day Wildfire Spread Prediction with a U-Net\\Phase 2: Deep Learning Report}

\author{\IEEEauthorblockN{Vanshika, Riya, Krish, Nitanshu, Tejas}
\IEEEauthorblockA{Team TensorTitans, Section C\\
Machine Learning Course Project}}

\maketitle

\begin{abstract}
We predict which 1~km squares will burn the next day using the Next Day Wildfire Spread dataset. Our Phase 1 tabular models learned nothing: with only 62 training fire pixels they scored a PR-AUC of 0.011, equal to random. In Phase 2 we train a small U-Net (491{,}553 parameters) on all 14{,}979 training maps (658{,}270 fire pixels), with wind as two vectors, masked unknown pixels, a weighted loss, flip augmentation that also flips the wind, dropout, weight decay and early stopping. On 6.7 million held-out test pixels the U-Net reaches a PR-AUC of 0.356, which is 28 times the random level and about twice the best simple rule, and it catches 56\% of fire pixels with 32\% precision. An ablation study shows that the fire mask of today matters most. All thresholds are chosen on validation data only.
\end{abstract}

\begin{IEEEkeywords}
wildfire spread, U-Net, semantic segmentation, class imbalance, precision-recall
\end{IEEEkeywords}

\section{Introduction}
Wildfires burn forests, homes and air quality, and a fire can change a lot within a single day. If a fire manager knows which places are likely to burn \emph{tomorrow}, they can send crews and warn people earlier. This project predicts, for every 1~km square on a map, whether it will be on fire the next day.

We use the Next Day Wildfire Spread (NDWS) dataset \cite{huot2022nextday}. Each sample is a $64\times64$ map with layers such as wind, temperature, humidity, drought, vegetation, terrain and the fire mask of today. The label is the fire mask of tomorrow. Fire is rare: only about 1.1\% of known pixels burn.

\subsection{What Phase 1 taught us}
In Phase 1 we turned each pixel into a row of a table and trained tree models. The result was honest but poor. Because of limited memory we used only 30 maps, which contained 62 fire pixels. On the held-out test the best model (LightGBM) reached a PR-AUC of 0.011, which is the same as a random guess, and it caught none of the 157 fire pixels in its sampled test set. Two simple rules, ``tomorrow looks like today'' and ``count the burning neighbours'', did better (0.109 and 0.182 on that test set), so the signal was in the data but the tabular models could not use it. Two reasons stand out: a pixel-by-pixel table cannot see the shape of a fire, and 62 fire pixels are too few to learn from.

\subsection{What Phase 2 does}
Deep learning fixes both problems. A convolutional network reads the whole map, so it sees neighbours and shapes, and it can be trained on every map without building a table. Our contributions are:
\begin{itemize}
\item A data pipeline that keeps all 18{,}545 maps (658{,}270 training fire pixels instead of 62), builds wind as two vectors, ignores unknown pixels and fits normalisation on the training split only.
\item A small U-Net with 491{,}553 parameters, trained on a laptop CPU in about 31 minutes, with a weighted loss, flip augmentation that also flips the wind, dropout, weight decay and early stopping.
\item A held-out evaluation on 6.7 million test pixels with a threshold chosen on validation. The U-Net reaches a test PR-AUC of 0.356, which is 28 times the random level and about twice the best simple rule.
\item An ablation study of seven design choices, each scored on validation only.
\end{itemize}

The rest of the report is organised as follows. Section~II reviews related work, Section~III describes the data, Section~IV the model and training, Section~V the results, Section~VI the ablation and Section~VII concludes.

\section{Related Work}
\subsection{The NDWS dataset and its baselines}
Huot et al.\ \cite{huot2022nextday} released the NDWS dataset and trained a convolutional autoencoder, logistic regression and a random forest. The autoencoder reached a test PR-AUC of about 0.284, the persistence rule about 0.115, and the two non-deep models were close to each other and clearly below the network. We use this paper as our reference. Their training pipeline is not identical to ours, so we compare scores as ``comparable'' rather than claiming a win. Fitzgerald et al.\ \cite{fitzgerald2023paying} built a U-Net with attention blocks for the same data, which confirms that U-Net style models are a sensible choice here.

\subsection{U-Net and its variants}
The U-Net \cite{ronneberger2015unet} has an encoder that finds what is in the image, a decoder that finds where, and skip connections that carry fine detail across. It was designed for medical images and works with little data and heavy augmentation. Attention U-Net \cite{oktay2018attention} adds gates on the skip connections so that the decoder focuses on useful regions. We use the plain U-Net and list the attention version as future work.

\subsection{Losses for rare positives}
When almost every pixel is negative, plain cross-entropy is dominated by easy negatives. Focal loss \cite{lin2017focal} reduces the weight of easy pixels, and the Dice loss \cite{milletari2016vnet} rewards overlap with the true shape. We use weighted cross-entropy and test focal loss in our ablation, where the weighted version was better.

\subsection{Evaluation under class imbalance}
Saito and Rehmsmeier \cite{saito2015precision} show that the precision-recall curve is more informative than the ROC curve when positives are rare, which is why our main metric is PR-AUC. Kumar et al.\ \cite{kumar2025stochastic} argue that matching fire pixels is not the same as learning the random process behind fire and suggest checking calibration. We did not measure calibration and list it as future work.

\subsection{Time series and other model families}
Models with memory, such as ConvLSTM \cite{shi2015convlstm}, need a sequence of days. NDWS gives only one input day, whereas WildfireSpreadTS \cite{gerard2023wildfirespreadts} provides full daily series for 607 fires, which is the right setting for such models. Vision transformers and Swin U-Net \cite{dosovitskiy2021vit,cao2022swinunet} usually need much more data and a GPU than a $64\times64$ map and a laptop deadline allow. Stacking a second model on top of the network \cite{wolpert1992stacked} needs predictions on data the network has not seen, which makes it expensive to do correctly; we describe our early attempt in the conclusion.

\section{Dataset and Preprocessing}
\subsection{Data and splits}
NDWS contains 18{,}545 maps of size $64\times64$ at 1~km per pixel, covering the contiguous United States from 2012 to 2020. The authors provide the train, validation and test split as separate TFRecord files, and we keep it unchanged (Table~\ref{tab:splits}). In Phase 1 only 30 maps were used. Here every map of every split is used. Unknown pixels (label $-1$, caused by clouds or missing data) are not counted as ``no fire''; they get weight zero in the loss and are left out of every metric.

\begin{table}[t]
\centering
\caption{Maps and fire pixels in each split.}
\label{tab:splits}
\footnotesize
\begin{tabular}{lrrrr}
\toprule
Split & Maps & Fire pixels & No-fire pixels & Fire rate \\
\midrule
Train & 14{,}979 & 658{,}270 & 59{,}254{,}803 & 1.10\% \\
Validation & 1{,}877 & 102{,}594 & 7{,}397{,}791 & 1.37\% \\
Test & 1{,}689 & 84{,}331 & 6{,}643{,}368 & 1.25\% \\
\bottomrule
\end{tabular}
\end{table}

\subsection{Input channels}
There are 12 raw input layers: elevation, wind direction, wind speed, minimum and maximum temperature, specific humidity, precipitation, the Palmer drought index, vegetation (NDVI), population density, the energy release component and the fire mask of today. Wind direction is an angle, so $359^\circ$ and $1^\circ$ look far apart as numbers although they are nearly the same direction. The angle also tells where the wind comes \emph{from}. We therefore replace it with two vectors that point where the wind blows \emph{to}:
\begin{equation}
w_x = -v\sin\theta, \qquad w_y = -v\cos\theta,
\end{equation}
where $v$ is wind speed and $\theta$ the direction. Wind speed is kept, which gives 13 input channels. Unknown values in today's fire mask are set to 0. Figure~\ref{fig:sample} shows one map: the arrows point towards the east and south-east, which is where the fire grew the next day.

\begin{figure*}[t]
\centering
\includegraphics[width=\textwidth]{dl_sample_map.png}
\caption{One training map: elevation, maximum temperature, fire today with wind arrows (blowing to), and fire tomorrow (target).}
\label{fig:sample}
\end{figure*}

\subsection{Normalisation without leakage}
Every channel is clipped to its 0.1 to 99.9 percentile range, which removes extreme values, and then standardised to mean 0 and standard deviation 1. All clip limits, means and standard deviations are computed from the training split only and reused for validation and test. The fire mask stays 0 or 1.

\subsection{Augmentation and regularisation}
During training, a mini-batch is flipped left-right with probability 0.5 and up-down with probability 0.5. Flipping the map changes the direction of the wind, so a left-right flip also changes the sign of $w_x$ and an up-down flip changes the sign of $w_y$ (Figure~\ref{fig:flip}). Without this, the wind would point the wrong way in the flipped map. We do not rotate maps. Other regularisers are described in Section~IV.

\begin{figure}[t]
\centering
\includegraphics[width=\columnwidth]{dl_wind_flip_check.png}
\caption{A left-right flip changes the sign of $w_x$, so the wind still points towards the same side of the fire.}
\label{fig:flip}
\end{figure}

\section{Methodology}
\subsection{U-Net architecture}
The network (Figure~\ref{fig:unet}) takes a $64\times64\times13$ input and returns a $64\times64\times1$ map of fire probabilities. The encoder has three levels with 16, 32 and 64 filters. Each level applies two $3\times3$ convolutions, each followed by batch normalisation \cite{ioffe2015batchnorm} and ReLU, and then a $2\times2$ max pool, so the map shrinks from 64 to 32, 16 and 8 pixels. The bottleneck has 128 filters at $8\times8$ with dropout 0.2 \cite{srivastava2014dropout}. The decoder upsamples back to $64\times64$; at every level the matching encoder output is joined to the upsampled map (skip connection), so edge detail is not lost. A final $1\times1$ convolution with a sigmoid gives the probability of fire at every pixel. The model has 491{,}553 parameters (1.88~MB).

\begin{figure*}[t]
\centering
\includegraphics[width=0.9\textwidth]{dl_unet_architecture.png}
\caption{The U-Net used in this project.}
\label{fig:unet}
\end{figure*}

\subsection{Loss}
Let $y_i\in\{0,1\}$ be the true label and $p_i$ the predicted probability of pixel $i$, and let $N$ be the number of pixels in a batch. The loss is a weighted binary cross-entropy:
\begin{equation}
\mathcal{L} = -\frac{1}{N}\sum_{i} w_i\left[y_i\log p_i + (1-y_i)\log(1-p_i)\right],
\end{equation}
with $w_i = 3$ for fire pixels, $w_i=1$ for no-fire pixels and $w_i=0$ for unknown pixels. The weight 3 makes a missed fire cost more than a false alarm, because fire is only about 1\% of pixels.

\subsection{Training setup}
We use the AdamW optimiser \cite{loshchilov2019adamw} with learning rate $10^{-3}$ and weight decay $10^{-4}$, mini-batches of 32 maps and at most 25 epochs. The learning rate is halved when validation PR-AUC does not improve for 2 epochs. Training stops when validation PR-AUC has not improved for 5 epochs and the weights of the best epoch are restored. The random seed is fixed to 42, and two runs with the same seed gave identical predictions in a short check (1{,}000 maps, one epoch). The code uses Keras \cite{chollet2015keras} and runs on the CPU of a laptop (Apple M3).

\subsection{Evaluation protocol}
Scores are computed on all pixels with a known label. Our main metric is PR-AUC (average precision),
\begin{equation}
\mathrm{AP}=\sum_n (R_n-R_{n-1})P_n,
\end{equation}
where $P_n$ and $R_n$ are the precision and recall at the $n$-th threshold. A random guess scores the fire rate, so PR-AUC is only meaningful when it is clearly above that rate. We also report ROC-AUC, recall at 5\% false-alarm rate, precision, recall and F1. The decision threshold is the one with the best F1 on the validation set and is frozen before the test set is used. The test set is used once, for the final model. Two simple rules are scored on exactly the same pixels as baselines: ``burning today'' (the fire mask of today as the score) and ``count burning neighbours'' (the number of burning pixels in the $3\times3$ box around the pixel).

\section{Results}
\subsection{Training}
The U-Net trained for 18 epochs in about 31 minutes, and early stopping restored epoch 13, which had a validation PR-AUC of 0.298 (Figure~\ref{fig:curves}). Training PR-AUC kept rising (0.64 at epoch 18) while validation stayed near 0.29 to 0.30, which is mild overfitting; early stopping kept the best epoch. Training PR-AUC is also higher partly because fire pixels count three times there. Validation loss stays below training loss because dropout and augmentation are active only during training.

\begin{figure}[t]
\centering
\includegraphics[width=\columnwidth]{dl_training_curves.png}
\caption{Loss and PR-AUC per epoch. The dashed line marks the epoch kept by early stopping.}
\label{fig:curves}
\end{figure}

\subsection{Held-out test performance}
Table~\ref{tab:test} compares the U-Net with random guessing and the two simple rules on the same 6{,}727{,}699 test pixels, using thresholds chosen on validation. The test fire rate is 1.25\%, so a random guess scores a PR-AUC of 0.0125. The U-Net reaches 0.356, which is 28 times the random level and about twice the best rule (0.173). Its ROC-AUC is 0.944. With the threshold 0.343 it catches 47{,}338 of the 84{,}331 fire pixels (recall 56\%) with a precision of 32\%. At a 5\% false-alarm rate it catches 79\% of fire pixels, compared with 49\% for the neighbour rule. The Phase 1 LightGBM scored 0.011 (ROC-AUC 0.495) on its own sampled test set of 14{,}429 pixels, so this row is a reference and not a same-pixel comparison.

\begin{table}[t]
\centering
\caption{Held-out test results (6.7 million pixels).}
\label{tab:test}
\footnotesize
\begin{tabular}{lccccc}
\toprule
Model & PR-AUC & ROC-AUC & Recall & Prec. & F1 \\
\midrule
Random guess & 0.013 & 0.500 & -- & -- & -- \\
Rule: burning today & 0.107 & 0.633 & 0.27 & 0.36 & 0.31 \\
Rule: burning neighbours & 0.173 & 0.729 & 0.47 & 0.25 & 0.33 \\
Phase 1 LightGBM$^{*}$ & 0.011 & 0.495 & 0.00 & 0.00 & 0.00 \\
\textbf{U-Net} & \textbf{0.356} & \textbf{0.944} & \textbf{0.56} & \textbf{0.32} & \textbf{0.40} \\
\bottomrule
\multicolumn{6}{l}{\scriptsize $^{*}$Phase 1 test set (sampled), not the same pixels.}
\end{tabular}
\end{table}

\begin{figure}[t]
\centering
\includegraphics[width=\columnwidth]{dl_test_pr_and_confusion.png}
\caption{Test precision-recall curves (left) and the U-Net confusion matrix (right).}
\label{fig:prcm}
\end{figure}

\subsection{The accuracy trap}
The U-Net has an accuracy of 97.93\%, but always predicting ``no fire'' scores 98.75\%. Accuracy therefore says nothing here, and we rely on PR-AUC and F1. The confusion matrix in Figure~\ref{fig:prcm} shows 6{,}540{,}776 correctly quiet pixels, 102{,}592 false alarms, 36{,}993 missed fire pixels and 47{,}338 caught fire pixels.

\subsection{Example maps}
Figure~\ref{fig:examples} shows test maps chosen by fire size (90th, 50th and 10th percentile among maps with fire tomorrow), not by hand. The U-Net places probability on the right areas and follows the direction of growth, but its output is smooth. It cannot draw the ragged edge of a real fire, which explains why precision is about 32\%.

\begin{figure}[t]
\centering
\includegraphics[width=0.85\columnwidth]{dl_test_example_maps.png}
\caption{Fire today, true fire tomorrow and U-Net probability for three test maps.}
\label{fig:examples}
\end{figure}

\subsection{Comparison with the dataset paper}
The dataset paper reports a test PR-AUC of about 0.284 for its convolutional autoencoder \cite{huot2022nextday}, and our persistence rule scores 0.107 against their 0.115. Our U-Net scores 0.356. Because the training and evaluation pipelines are not identical, we call our result comparable or better and not a direct win.

\section{Ablation Study}
To see which design choices matter, we changed one thing at a time and measured validation PR-AUC. Training the full model takes about 31 minutes, so every run uses the same smaller budget: the first 4{,}000 training maps, at most 10 epochs, early stopping with patience 3 and the same seed. The reference run therefore scores 0.260, lower than the full model (0.298). The test set was not used. Table~\ref{tab:ablation} and Figure~\ref{fig:ablation} show the results.

\begin{table}[t]
\centering
\caption{Ablation results (validation PR-AUC).}
\label{tab:ablation}
\footnotesize
\begin{tabular}{lcc}
\toprule
Experiment & PR-AUC & Change \\
\midrule
Full U-Net (reference) & 0.260 & -- \\
No augmentation & 0.236 & $-0.023$ \\
Flips without wind fix & 0.255 & $-0.004$ \\
No fire class weight & 0.233 & $-0.026$ \\
Focal loss instead of weighted BCE & 0.231 & $-0.028$ \\
Without fire today & 0.175 & $-0.084$ \\
Without weather channels & 0.243 & $-0.017$ \\
\bottomrule
\end{tabular}
\end{table}

\begin{figure}[t]
\centering
\includegraphics[width=\columnwidth]{dl_ablation.png}
\caption{Ablation: validation PR-AUC when one component is changed.}
\label{fig:ablation}
\end{figure}

\textbf{Fire today matters most.} Removing the fire mask of today drops PR-AUC by 0.084. It stays above random because the maps are centred on fires, so the other channels still hint where the fire is.

\textbf{Loss.} Weighting fire pixels three times beats plain cross-entropy ($-0.026$) and focal loss ($-0.028$). Focal loss was designed for this kind of imbalance, but on our data the simple weight worked better.

\textbf{Augmentation and weather.} Flips help ($-0.023$ without them), and weather channels help a little ($-0.017$ without them).

\textbf{Wind fix.} Flipping without changing the wind is only 0.004 worse. This is too small to call a real effect from one seed, but changing the wind is still the physically correct thing to do.

\textbf{Limits.} Each row is one run with one seed, and the differences between the smaller changes may be within random variation. The ranking of the large effects, especially the fire mask of today, is the reliable part.

\section{Conclusion}
We moved from a Phase 1 table model that learned nothing (test PR-AUC 0.011, equal to random, 62 training fire pixels) to a U-Net that reads whole maps and trains on all 14{,}979 maps (658{,}270 fire pixels). On a held-out test of 6.7 million pixels it reaches a PR-AUC of 0.356, 28 times the random level and about twice the best simple rule, and it catches 56\% of fire pixels at a precision of 32\%. The ablation shows that the fire mask of today matters most, and that class weighting, flip augmentation and weather channels each help.

\subsection{Limitations}
\begin{itemize}
\item Every result comes from one training run with one seed, so we cannot give confidence intervals.
\item The ablation uses a smaller budget than the final model, so only its ranking is meaningful.
\item The decision threshold and early stopping use the same validation set.
\item The comparison with the dataset paper is approximate because the pipelines differ.
\item The output is smooth, so edge precision is limited (precision about 32\%).
\item We did not check calibration, which \cite{kumar2025stochastic} recommends.
\end{itemize}

\subsection{Future work}
The natural next steps are an Attention U-Net \cite{oktay2018attention}, a Dice-based loss \cite{milletari2016vnet}, several seeds for confidence intervals, a calibration check, and a proper stacked hybrid with a gradient-boosted model. We tried a quick draft of the hybrid before the final notebooks: it scored higher on validation but lower on test than the U-Net alone. Because the second model had only half of the validation maps to learn from and the script is not part of our final code, we do not claim it. A correct version needs out-of-fold predictions \cite{wolpert1992stacked}. Finally, datasets with several input days \cite{gerard2023wildfirespreadts} would make models with memory possible.


\bibliographystyle{IEEEtran}
\bibliography{references}

\end{document}
